% Figure captions. Generated by scripts/figures_v3/build.py and edited for the
% delivery package: implementation vocabulary replaced by standard terminology, and
% every caption states whether it displays a proved bound, an exact construction, or an
% observed value.
\newcommand{\FigCaptionOne}{Semantics of a typed ordered composition tree. Types
determine the ambient spaces, the isometric embeddings, the orthogonal projectors and
the admissible multilinear maps; unprojected evaluation, recursive projection and the
local closure residual are three distinct operations. This is a definitional diagram: it
carries no sampling uncertainty and supports no statement about optimal constants.}

\newcommand{\FigCaptionTwo}{Exact local error decomposition. The local closure residual
and every nonempty subset of propagated child errors appear in the multilinear identity,
followed by the orthogonal splitting into projected and orthogonal components. This is a
diagram of a proof; it is an identity, not a bound, and it does not determine optimal
constants.}

\newcommand{\FigCaptionThree}{Comparison of four upper bounds on 4{,}320 distinct
configurations with heterogeneous local parameters (10 declared parameter draws; no
optimiser restarts are involved, since all four bounds are computed in closed form).
Each bound is divided by the uniform ambient bound of Theorem~\ref{thm:homogeneous} on
the same configuration; points are medians and bars are the 10th--90th percentile range
over configurations, not an uncertainty estimate. The data show that the vertexwise,
pathwise and order-optimised bounds improve on the uniform bound when the local
parameters are nonuniform. They do not show that any of these bounds is optimal.}

\newcommand{\FigCaptionFour}{Ratios of the best explicit lower construction to the
proved upper bound, over 240 two-dimensional configurations selected from 2{,}880. Tile
colour is the ratio; the text gives the enclosed lower value, the upper bound, and the
topology attaining the maximum. Lower values are rigorous interval enclosures of an
explicit construction, so they carry no sampling uncertainty. A ratio of one identifies
a configuration in which the constant is determined; ratios below one leave the exact
constant undetermined.}

\newcommand{\FigCaptionFive}{Explicit lower constructions for the projected-error
constant on 96 distinct binary trees with eight internal vertices. The vertical axis is
the interval-enclosed projected error divided by $\eta M^{k-1}$; the dashed line is the
proved upper coefficient $k-1=7$. The construction is explicit, so there is no
statistical uncertainty. The figure exhibits an admissible lower profile; it does not
determine $C_T^{\mathrm{proj}}(\eta)$.}

\newcommand{\FigCaptionSix}{The coefficients $k$ and $k-1$ on 10 binary left-comb trees
in dimension two at the smallest tested $\eta$, with interval-enclosed lower
constructions. Ambient and projected errors are normalised by $\eta M^{k-1}$ and
compared with the proved coefficients. The ambient construction attains the linear
coefficient in these configurations. The projected construction shown here, at small
$\eta$, remains far below $k-1$; whether $k-1$ is attained at fixed $\eta>0$ is open.}

\newcommand{\FigCaptionSeven}{Dependence of the bound on tree topology, over 240
distinct configurations at $\eta=0.1$. The axes are tree depth and the sum of
root-to-vertex path lengths; colour encodes imbalance, marker size the Strahler number,
and a highlighted edge marks a projected lower-to-upper ratio above $0.9$. Topology
statistics are exact and lower values are interval-enclosed. The figure shows that the
bound depends on topology; it does not identify a transition or an extremal topology.}

\newcommand{\FigCaptionEight}{Attribution of the projected root error to individual
vertices, for one balanced tree with eight internal vertices. All $2^{8}=256$ subsets of
vertices at which the closure residual is switched on were evaluated exhaustively; no
sampling is involved. Vertex labels give the Shapley value of the resulting set
function, and the lower panel compares single-vertex main effects with the full Shapley
values. The efficiency residual is below $10^{-12}$. The attribution is exact for this
construction and is not a claim about a unique causal mechanism.}

\newcommand{\FigCaptionNine}{Interaction structure of the local error sources, from all
1{,}530 exhaustive subset evaluations over 24 trees. For each tree, the exact set
function of the projected error is M\"obius-inverted; bars give the topology-median
fraction of absolute coefficient mass carried by first-, second- and higher-order
subsets. The reference is the exact subset expansion of Theorem~\ref{thm:subset}. The
data show that the local error sources do not combine additively in these
constructions.}

\newcommand{\FigCaptionTen}{Effect of the ordering rule of Theorem~\ref{thm:order}, over
all 4{,}320 distinct configurations. The vertical axis is the bound obtained with the
optimal order divided by the bound obtained with the reference order. Boxes give the
median, quartiles and 5th--95th percentile range over configurations. The rule is never
worse and is sometimes strictly better. It is optimal within the scalar telescoping
family of Section~\ref{sec:order} only, not among all bounding strategies.}

\newcommand{\FigCaptionEleven}{Geometry of the extremal construction. Panels (a) and (b)
display the projected and orthogonal output blocks of one explicit rank-one gated
rotation tensor; panel (c) summarises the mechanism by which an orthogonal error created
at a lower vertex acquires a projected component at its parent. Entries are the exact
stored inputs up to floating-point representation. The figure establishes admissibility
and alignment of this construction, not its uniqueness or global optimality.}

\newcommand{\FigCaptionTwelve}{Enclosures of the constant for 36 small binary
configurations at $\eta=0.1$. Each horizontal interval joins the interval-enclosed lower
construction to the proved upper bound; a filled marker identifies a configuration in
which the two coincide. The references are the coefficients $k$ and $k-1$. Only the
marked configurations have a determined constant; the remaining intervals are open, and
their width is not an error estimate.}

\newcommand{\FigCaptionThirteen}{Signed linear combinations of composition trees: seven
families over 28 configurations. The exact quantity plotted is the normalised residual
of the signed combination. The two upper bars are proved upper bounds --- the triangle
bound and the bound after exact combination of syntactically identical trees --- while
the third bar is the largest lower value found by numerical search over $\eta$. The
comparison shows that exact combination improves the bookkeeping for these declared
expressions. The search values are lower bounds and do not determine the constants for
the associator, the Filippov identity, or the generalised Jacobi expressions.}

\newcommand{\FigCaptionFourteen}{Separation of representation error from
recursive-projection error, for five configurations in dimension 64 on a balanced tree
with eight internal vertices at $\eta=0.01$. The three bars are the separately proved
representation, closure and interaction terms of
Proposition~\ref{prop:approximation}. Their sum is by definition the proved upper bound
and is not asserted to equal the observed error. The figure displays a decomposition of
a bound; it says nothing about tightness or about identifiability of the low-rank
representation.}

\newcommand{\FigCaptionFifteen}{Dependence on ambient dimension and projector rank, over
72 configurations on binary trees with 24 internal vertices at $\eta=0.1$. Each tile is
the value of the two-dimensional construction, embedded isometrically into the stated
dimension, divided by the proved upper coefficient $k-1$. Values are interval-enclosed.
The embedding preserves the value, so the construction is admissible in every dimension
shown. No configuration exhibits an improvement obtained from the additional
coordinates.}

\newcommand{\FigCaptionSixteen}{Convergence of the numerical search on a single
configuration: four gradient-based trajectories and one independent derivative-free
run. The vertical axis is the attained normalised lower value divided by the proved
upper bound. Every recorded step of every run is shown. Optimiser restarts are not
statistical replicates and no interval is inferred from them. Agreement close to one
indicates a strong lower construction for this configuration; it is not a certificate of
global optimality.}

\newcommand{\FigCaptionSeventeen}{Wall-time benchmark on the hardware described in
Appendix~\ref{app:reproducibility}: 15 backend and problem-size configurations, five
timing repetitions each, 15--30 evaluations per repetition. Panels show median wall time
with the observed minimum--maximum range, derived vertex throughput, and peak allocated
device memory, for the reference coordinate evaluator, the vectorised CPU evaluator and
the CUDA evaluator. Double-precision outputs agree exactly across the three. This is a
measurement on one machine: it supports no asymptotic-complexity statement and no
comparison between hardware platforms.}

\newcommand{\FigCaptionEighteen}{Dependency graph relating the statements of this
article to their supporting material. Proved upper bounds, rigorous lower enclosures,
exploratory numerical evidence, and unresolved questions are kept in separate layers.
The graph is a navigation aid for the reproducibility appendix; it does not convert a
numerical search result into a theorem.}
